\chapter{多重积分}
\label{chap:multiple-integrals}

一元 Riemann 积分把区间分成小段；在 $\R^n$ 中，小段由小矩形取代。
定义本身仍由 Riemann 和给出，但一般定义域的边界可能占据不可忽略的体积。
因此本章先在矩形上建立积分，再用 Jordan 零测边界说明哪些有界集合可以由
有限矩形稳定逼近。累次积分和变量代换另列下一章；那里需要这里建立的集合
可测性与有限可加性。

本章的矩形均为有界闭矩形
$Q=\prod_{j=1}^n[a_j,b_j]$，并记
$|Q|=\prod_j(b_j-a_j)$。退化矩形的体积约定为零。

\section{矩形上的分割与 Riemann 和}

在每个坐标区间选取有限分割，所有坐标小区间的笛卡儿积组成 $Q$ 的一个
矩形分割 $P$。以子矩形 $R$ 的直径最大值为网格 $\norm P$。在每个
非退化 $R\in P$ 中取标记点 $\xi_R$，定义
\[
 S(f,\dot P)=\sum_{R\in P}f(\xi_R)|R|.
\]

\begin{definition}
有界函数 $f:Q\to\R$ 称为 Riemann 可积，如果存在 $I\in\R$，使对
任意 $\varepsilon>0$，存在 $\delta>0$，只要 $\dot P$ 是标记矩形分割且
$\norm P<\delta$，就有
\[
 |S(f,\dot P)-I|<\varepsilon.
\]
此时写作 $I=\int_Q f(x)\,dx$。
\end{definition}

\section{Darboux 判据与连续函数}

令 $M_R=\sup_Rf$、$m_R=\inf_Rf$，并定义
\[
 U(f,P)=\sum_RM_R|R|,\qquad L(f,P)=\sum_Rm_R|R|.
\]
一元情形的公共加细论证逐坐标照搬，得到加细使上和不增、下和不减。
于是可以定义
\[
 \overline{\int_Q}f=\inf_PU(f,P),\qquad
 \underline{\int_Q}f=\sup_PL(f,P).
\]
任意两个分割都有公共加细，所以每个下和不超过每个上和，从而
$\underline{\int_Q}f\le\overline{\int_Q}f$。二者相等时称 $f$ Darboux
可积，其公共值称为 Darboux 积分。

\begin{theorem}[Darboux 判据]
\label{thm:multiple-darboux-criterion}
有界函数 $f:Q\to\R$ 可积，当且仅当对每个 $\varepsilon>0$，存在矩形
分割 $P$ 使
\[
 U(f,P)-L(f,P)=\sum_{R\in P}\omega(f;R)|R|<\varepsilon.
\]
Riemann 定义与 Darboux 定义给出同一积分。
\end{theorem}

\begin{proof}
先设振幅判据成立，并记 $|f|\le M$。由
$\underline{\int_Q}f\ge L(f,P)$ 与
$\overline{\int_Q}f\le U(f,P)$，
\[
 0\le \overline{\int_Q}f-\underline{\int_Q}f
 \le U(f,P)-L(f,P).
\]
右端可以任意小，故上下积分相等，记公共值为 $I$。给定
$\varepsilon>0$，选分割 $P_0$
使 $U(f,P_0)-L(f,P_0)<\varepsilon/3$。$P_0$ 的内部分割面只有有限多个；
存在只依赖于 $P_0,Q$ 的常数 $C$，使任一网格为 $\eta$ 的矩形分割中，
所有与这些分割面相交的小矩形总体积不超过 $C\eta$。取 $\eta$ 充分小，
使 $2MC\eta<\varepsilon/3$。

任取网格小于 $\eta$ 的标记分割。除上述小矩形外，每个小矩形都包含在
$P_0$ 的某一块内，其标记值介于该块的下、上确界之间。坏矩形上任意两个
可能取值的差至多为 $2M$。因此标记和 $S(f,\dot P)$ 与区间
$[L(f,P_0),U(f,P_0)]$ 的距离小于 $\varepsilon/3$。上下积分的公共值
$I$ 也在这个区间内，于是
\[
 |S(f,\dot P)-I|
 \le U(f,P_0)-L(f,P_0)+\frac{\varepsilon}{3}
 <\frac{2\varepsilon}{3}.
\]
这证明所有充分细的标记和都趋于 $I$。

反过来，设 Riemann 定义成立。给定 $\varepsilon>0$，取 $\delta>0$，
使网格小于 $\delta$ 的任一标记和与积分 $I$ 的距离小于
$\varepsilon/4$，并固定一个这样的分割 $P$。由于 $P$ 只有有限块，可在
每块分别选两个标记点，使对应函数值距该块的上、下确界都小于一个共同的
充分小正数。所得两组标记和 $S^+,S^-$ 满足
\[
 U(f,P)-S^+<\frac{\varepsilon}{4},\qquad
 S^--L(f,P)<\frac{\varepsilon}{4}.
\]
又有 $|S^\pm-I|<\varepsilon/4$，故
$U(f,P)-L(f,P)<\varepsilon$。振幅判据成立，两种定义的积分值相同。
\end{proof}

\begin{theorem}[连续函数可积]
\label{thm:continuous-multiple-riemann-integrable}
若 $f:Q\to\R$ 连续，则 $f$ 在 $Q$ 上 Riemann 可积。
\end{theorem}

\begin{proof}
若 $|Q|=0$，按约定积分为零。以下设 $|Q|>0$。$Q$ 紧，故 $f$ 一致连续。
给定 $\varepsilon>0$，取 $\delta>0$，使
$|x-y|<\delta$ 时 $|f(x)-f(y)|<\varepsilon/|Q|$。网格小于
$\delta$ 的分割满足
\[
 U(f,P)-L(f,P)
 \le \frac{\varepsilon}{|Q|}\sum_{R\in P}|R|=\varepsilon.
\]
\end{proof}

线性、序关系、乘积和绝对值的可积性均由振幅估计得到；特别地
\[
 \left|\int_Qf\right|\le\int_Q|f|\le |Q|\sup_Q|f|.
\]

\section{Jordan 零测集}

\begin{definition}
有界集 $N\subset\R^n$ 称为 Jordan 零测集，如果对每个
$\varepsilon>0$，存在有限多个矩形 $R_1,\ldots,R_m$ 覆盖 $N$，且
$\sum_k|R_k|<\varepsilon$。有界集 $E$ 称为 Jordan 可测，如果其边界
$\partial E$ 是 Jordan 零测集。
\end{definition}

有限集、光滑曲线在 $\R^2$ 中的像以及光滑曲面在 $\R^3$ 中的像都是
Jordan 零测集。以曲线为例，一致连续性把参数区间分成有限小段，使每段像
落在边长受控的小正方形中；若曲线为 $C^1$，有界导数还给出总面积
$O(1/N)$ 的直接估计。稠密可数集则未必 Jordan 零测；Jordan 覆盖只允许
有限个矩形，因而 $\Q^n\cap Q$ 的闭包是整个 $Q$。

连续性是充分条件，却远非必要条件。多元积分中常见的可积函数会沿有限条
曲线或曲面发生跳跃；这些集合虽然包含无穷多个点，却不占据体积。

\begin{theorem}[Jordan 零测不连续集判据]
\label{thm:jordan-null-discontinuities-integrable}
设 $f:Q\to\R$ 有界，并设其全部不连续点包含在一个 Jordan 零测集 $N$ 中。
则 $f$ 在 $Q$ 上 Riemann 可积。
\end{theorem}

\begin{proof}
设 $|f|\le M$。若 $M=0$ 或 $|Q|=0$，结论直接成立。以下设 $M>0$ 且
$|Q|>0$。给定 $\varepsilon>0$，用有限个开矩形覆盖 $N$，使这些矩形
的体积和小于 $\varepsilon/(8M)$。把矩形的端点
加入坐标分割，并把网格进一步取细，可使所有与该覆盖相交的小矩形的总体积
小于 $\varepsilon/(4M)$。这些小矩形上的振幅至多为 $2M$，故它们对振幅和
的贡献小于 $\varepsilon/2$。

删去覆盖后留下 $Q$ 的一个紧子集 $K$，$f|_K$ 连续，因而一致连续。再把
网格取得足够细，使每个完全落在 $K$ 中的小矩形上的振幅小于
$\varepsilon/(2|Q|)$。其余小矩形已经计入前一类，于是
\[
 U(f,P)-L(f,P)<\frac{\varepsilon}{2}
 +\frac{\varepsilon}{2|Q|}\sum_{R\in P}|R|\le\varepsilon.
\]
Darboux 判据给出可积性。
\end{proof}

例如，若 $Q$ 被有限个 $C^1$ 曲面分成若干块，而函数在每一块的闭包上连续，
则它的不连续点至多落在这些分界曲面上，因而可积。这里要求的是对体积的
零测控制；仅仅说不连续点“很少”或“可数”并不足以适用于 Jordan 理论，
因为稠密可数集不能由总体积任意小的有限矩形族覆盖。

\section{Jordan 可测集}

\begin{theorem}[集合判据]
\label{thm:jordan-indicator-boundary}
设 $E$ 有界，$Q$ 是包含 $E$ 的矩形。以下条件等价：
\begin{enumerate}[label=\textup{(\roman*)}]
 \item $E$ Jordan 可测；
 \item 指示函数 $\mathbf1_E$ 在 $Q$ 上 Riemann 可积；
 \item 对每个 $\varepsilon>0$，存在有限个内部两两不交的矩形，使其并的
 对称差与 $E$ 包含于总体积小于 $\varepsilon$ 的有限矩形并中。
\end{enumerate}
积分 $\int_Q\mathbf1_E$ 与外包矩形无关，称为 $E$ 的 Jordan 体积
$|E|$。
\end{theorem}

\begin{proof}
在一个分割小矩形 $R$ 上，$\mathbf1_E$ 的振幅为一，当且仅当
$R$ 同时遇到 $E$ 和其补集；这类矩形必与 $\partial E$ 相交。若边界可由
总体积任意小的有限矩形覆盖，把覆盖端点加入分割，振幅判据给出可积性。

反之，若指示函数可积，取振幅和小于 $\varepsilon$ 的分割。所有振幅为一
的小矩形覆盖 $\partial E$，且总体积小于 $\varepsilon$，故边界为零测集。
 把振幅为零且函数恒为一的小矩形并起来，便得到 (iii)。

 设反过来 (iii) 成立，记这些内部两两不交矩形的并为 $R$，并设
 $E\mathbin\triangle R$ 包含于有限个闭矩形之并 $G$，其总体积小于
 $\varepsilon$。集合的边界满足
 \[
  \partial E\subseteq \partial R\cup\overline{E\mathbin\triangle R}
  \subseteq\partial R\cup G.
 \]
 有限矩形并的边界包含在有限多个矩形面的并中；把每个面置于任意薄的矩形
 条内，说明 $\partial R$ 为 Jordan 零测集。因此，另取总体积小于
 $\varepsilon$ 的有限矩形覆盖 $\partial R$，便得到总体积小于
 $2\varepsilon$ 的 $\partial E$ 覆盖。令 $\varepsilon$ 任意即得 (i)。

 外包矩形改变只会增加指标函数恒为零的区域，所以积分不变。
\end{proof}

若 $E$ Jordan 可测，$f:E\to\R$ 有界，并能在某个外包矩形上延拓为
$\widetilde f$，使 $\widetilde f\mathbf1_E$ 可积，则定义
\[
 \int_E f=\int_Q\widetilde f\mathbf1_E.
\]
对 $\overline E$ 上连续的 $f$，Tietze 延拓定理或在边界小矩形上的直接
振幅估计保证定义成立，且与延拓无关。

\section*{习题}
\addcontentsline{toc}{section}{习题}

\subsection*{A组　定义与可积性}
\begin{enumerate}[label=\textbf{\thechapter.\arabic*.}]
 \exerciseitem{multiple-integral-product-function}
 由 Riemann 和证明 $f(x)g(y)$ 在矩形上的积分等于两个一元积分之积。
 \exerciseitem{multiple-continuous-integrability}
 补全连续函数多重可积性的振幅证明，并处理退化矩形。
 \exerciseitem{jordan-boundary-indicator}
 详细证明定理 \ref{thm:jordan-indicator-boundary} 中 (i) 与 (ii) 的等价。
 \exerciseitem{graph-jordan-null}
 证明闭区间上 Lipschitz 函数的图像在 $\R^2$ 中是 Jordan 零测集。
 \exerciseitem{dense-countable-not-jordan-null}
 证明 $\Q^n\cap[0,1]^n$ 不是 Jordan 零测集，并说明它与“可数集测度为零”
 的 Lebesgue 结论并不矛盾。
 \exerciseitem{jordan-volume-finite-union}
 设有限个 Jordan 可测集的内部两两不交。证明任意两个集合的交都包含于
 它们的边界之并，并由此证明 Jordan 体积的有限可加性。
\end{enumerate}

\subsection*{B组　集合运算与逼近}
\begin{enumerate}[label=\textbf{\thechapter.\arabic*.},resume]
 \exerciseitem{multiple-darboux-cauchy}
 证明矩形上有界函数的 Darboux 判据，并由此推出标记 Riemann 和的 Cauchy
 判据。
 \exerciseitem{jordan-set-algebra}
 证明 Jordan 可测集对有限并、有限交和差封闭，并写出相应边界包含关系。
 \exerciseitem{jordan-product-volume}
 若 $E\subset\R^p$、$F\subset\R^q$ Jordan 可测，证明
 $E\times F$ Jordan 可测且 $|E\times F|=|E|\,|F|$。
 \exerciseitem{jordan-continuous-domain}
 设 $E$ Jordan 可测，$f$ 在 $\overline E$ 上连续。证明本章定义的
 $\int_Ef$ 与所取连续延拓及外包矩形无关。
\end{enumerate}
